\documentclass[UTF8,a4paper]{ctexart}
\usepackage{longtable,booktabs,array,hyperref}
\begin{document}
\section*{2027 国科大杭高院理论物理备考总纲}
当前科目为 101 政治、201 英语（一）、604 高等数学、811 量子力学。604、811 均为150分、180分钟闭卷；811 为五道计算题。目标为理论物理研究所代国科大杭州高等研究院招生方向（070201），正式目录、导师、人数和考场时段以官方后续公告和准考证为最终依据。

\section{执行原则}
量子力学已完成基础通读，后续以推导、综合和限时书写为主；604 从系统基础启动，但按已有基础快速推进。9--10 月专业课有效时间按 604:811 约 3:2 分配，11 月后按整卷数据调整。每次学习完成“短时回忆--新内容/例题--不看答案做题--订正并标记原因”闭环；只看懂不算完成。新内容在当天、约三天后、约十天后各进行一次不看笔记的回忆或重做。每天安排三至四个 60--90 分钟专注块，块间休息 10--20 分钟；每周半天休息或机动，疲劳时保留回忆和订正，不熬夜补时。

\section{604 高等数学范围}
604 官方大纲为微积分约80\%和线性代数约20\%。每项均要达到：不看笔记能选择方法、写关键步骤，并在限时练习中稳定完成。
\begin{itemize}
\item 函数、极限、连续：函数性质、左右极限、等价无穷小、夹逼和单调有界、间断点、闭区间性质。
\item 一元微分与积分：导数、微分、隐函数和参数方程、高阶导数；中值定理、Taylor、洛必达；不定/定积分、反常积分及应用。
\item 向量与空间解析几何：直线平面、距离夹角、二次曲面、空间曲线和投影。
\item 多元微分与积分：全微分、隐函数、梯度、极值；二三重积分、线面积分、Green/Gauss/Stokes、路径无关、散度旋度。
\item 级数与常微分方程：敛散、幂级数、Taylor、Fourier；一阶、降阶、常系数和 Euler 方程。
\item 线性代数：行列式、矩阵和秩、向量组与正交化、方程组、特征值特征向量和相似对角化。
\end{itemize}
按“概念与典型例题--基础到中等难度习题--同源旧卷筛选训练”推进。旧 601 卷只作相近风格训练；超出 604 官方范围的题标为拓展，不占核心时间。

\section{811 量子力学范围}
每周至少一次空白纸推导：不看讲义完成一个模型或定理的骨架，再核对缺口。每题按“物理条件--表象或哈密顿量--边界或对称条件--计算--物理解释”书写。
\begin{itemize}
\item 波函数、Schrödinger 方程与一维势场：概率解释、归一化、连续性方程、定态与测量；方势阱、势垒、共振、delta 势、谐振子。
\item 算符、图像与表象：厄米性、共同本征、测量、不确定度、守恒量、Ehrenfest、两种图像；Dirac 符号、幺正变换、占有数表象。
\item 中心力场、自旋与耦合：两体约化、氢原子；Pauli 矩阵、单三重态、总角动量和 CG 系数。
\item 电磁场与全同粒子：最小耦合、规范、Zeeman、Landau；交换对称性、玻色费米与 Pauli 原理。
\item 定态近似与量子跃迁：非简并和简并微扰、变分；突发、绝热、周期和有限时微扰、选择定则。
\end{itemize}

\section{16 周节点}
\begin{longtable}{p{.18\textwidth}p{.35\textwidth}p{.35\textwidth}}
\toprule 阶段 & 604 主线 & 811 主线\\\midrule
第1--2周\\8月31日--9月13日 & 函数极限连续、行列式矩阵、导数微分；各一次90分钟诊断。 & 诊断一维势、算符、谐振子；完成势场与振子。\\
第3--4周\\9月14日--9月27日 & 中值定理、Taylor、积分、空间直线平面。 & 测量守恒量图像；中心力场与氢原子。\\
第5--6周\\9月28日--10月11日 & 二次曲面、多元微分极值、秩与方程组。 & 矩阵表象；自旋、角动量、CG系数。\\
第7--8周\\10月12日--10月25日 & 二三重积分、线面积分、三大公式、特征值对角化。 & 电磁场、全同粒子、微扰与变分；每科一套180分钟基线卷。\\
第9--10周\\10月26日--11月8日 & 级数、Fourier、常微分方程。 & 量子跃迁；按势场、算符、角动量分类真题。\\
第11--12周\\11月9日--11月22日 & 混合训练、弱项修复、线代维护。 & 中心力场、微扰、跃迁真题；五题配速，每周一套整卷。\\
第13--14周\\11月23日--12月6日 & 保留卷模拟，只补高频弱项。 & 保留卷模拟，只补高频弱项；一次半日联考。\\
第15--16周\\12月7日--12月20日 & 轻量限时、公式回忆、易错步骤。 & 轻量限时、模型骨架、易错步骤；停止刷新资料。\\\bottomrule
\end{longtable}

\section{日程、复盘与模拟}
普通日：604 专注块75--90分钟，811 专注块60--90分钟，604习题或订正60--90分钟；英语约90分钟，政治约60分钟，11月逐步增加政治背诵输出。每周复盘45--60分钟，只记录有效时长、限时正确率与耗时、错误类型和下周最多两项修复任务。错误只归为知识缺口、建模、推导、计算、时间管理。连续两次低于60\%的模块，下周增加两个45分钟修复块；连续两次达到80\%且时间合格，转为每周维护。

九月真题按模块使用；十月下旬开始整卷。整卷均在答题纸完成，记录每题起止时间和错误标签。每科至少保留两套未完整做过的卷给11--12月。模拟先浏览全卷；停滞时留条件后跳题；最后预留10--15分钟核对符号、边界、归一化、量纲和答案位置。

\section{官方依据与核验}
报名开始前核对2027年专业目录中的培养单位、代招方向、导师、人数和科目；考前以准考证确认日期、考点和时段。官方来源：\href{https://admission.ucas.ac.cn/showarticle/Article/4c7e0e9f-2311-47a0-8f12-b0ec992078ac/7d50f7f5-0125-4410-93a3-2b2b6917f0dd}{2027年初试科目设置公告}；\href{https://admission.ucas.ac.cn/Content/Upload/2026/7/29.pdf}{604高等数学考试大纲}；\href{https://admission.ucas.ac.cn/Content/Upload/2025/7/811%20%E9%87%8F%E5%AD%90%E5%8A%9B%E5%AD%A6.pdf}{811量子力学考试大纲}。
\end{document}
